% !TeX program = xelatex
\documentclass[12pt]{article}
\usepackage{makecell}
\usepackage{hyphenat}
\usepackage{fontspec}
% Load the TeX Live font files directly so the document also builds on
% machines where the system font index does not expose "TeX Gyre Termes".
\setmainfont{texgyretermes-regular.otf}[
  BoldFont=texgyretermes-bold.otf,
  ItalicFont=texgyretermes-italic.otf,
  BoldItalicFont=texgyretermes-bolditalic.otf
]
\usepackage{polyglossia}
\setdefaultlanguage{vietnamese}\setotherlanguage{english}
\usepackage[table,dvipsnames]{xcolor} % Required for row coloring
\usepackage{booktabs} % Required for nicer horizontal lines
\usepackage{ragged2e}
\usepackage{caption}
\usepackage{subcaption}
\usepackage{float}
\usepackage{amsmath,amssymb,amsfonts}
\usepackage{graphicx}
\usepackage{adjustbox}
\usepackage{multicol}
\usepackage[inline,shortlabels]{enumitem}
\usepackage{setspace}
\usepackage{longtable}
\usepackage{bm}
\usepackage{titling}
\usepackage{tabularx}
% fontenc is not needed with XeLaTeX
\usepackage{mathrsfs}
% \usepackage[draft]{hyperref}
\usepackage{hyperref}
\usepackage{xparse}
\usepackage{fontawesome5}
\usepackage[most]{tcolorbox}
\usepackage{tikz}
\usepackage[framemethod=TikZ]{mdframed}
\usetikzlibrary{matrix, fit, shadows, positioning, arrows.meta, calc, shapes.geometric, backgrounds, shapes.symbols}
\usepackage{paracol}
\usepackage{siunitx}
\usepackage{multirow}
\usepackage{forest}

\usepackage[toc,page]{appendix}
\renewcommand{\appendixtocname}{Phụ lục}
\renewcommand{\appendixpagename}{Phụ lục}

\usepackage{xurl}
\setlength{\headheight}{14.5pt}
\setlength{\topmargin}{-.5in} \setlength{\textheight}{9.25in}
\setlength{\oddsidemargin}{-0.1in} \setlength{\textwidth}{6.8in}

\usepackage{tocloft}
\renewcommand{\cfttoctitlefont}{\Huge\bfseries}
\setlength{\cftaftertoctitleskip}{2em}
\setlength{\cftbeforesubsecskip}{7pt} % giãn cách subitem mục lục
\renewcommand{\cftsecleader}{\cftdotfill{\cftdotsep}}
\setlength{\cftsecnumwidth}{3.2em}
\renewcommand{\cftsecaftersnum}{\quad}
\renewcommand{\cftsubsecleader}{\cftdotfill{\cftdotsep}}
\setlength{\cftsubsecnumwidth}{4em}
\renewcommand{\cftsubsecaftersnum}{\quad}
\renewcommand{\cftsubsubsecleader}{\cftdotfill{\cftdotsep}}
\setlength   {\cftsubsubsecnumwidth}{5em}
\renewcommand{\cftsubsubsecaftersnum}{\quad}

\newcounter{mycounter}
\newcommand\showmycounter{\stepcounter{mycounter}\themycounter}
\usepackage{lipsum}
\newcommand\showlips{\stepcounter{mycounter}\lipsum[\value{mycounter}]}

%--------------------------------------------------------------------
\usepackage{framed}
\usepackage{fancyhdr}
\usepackage{upquote}
\usepackage{listings}


% Color --------------------------------------------------------------------
\usepackage{epigraph}
\usepackage{titlesec}

\definecolor{codebackground}{RGB}{250, 250, 250}
\definecolor{codecomment}{RGB}{34, 139, 34}
\definecolor{codekeyword}{RGB}{0, 0, 255}
\definecolor{codestring}{RGB}{255, 0, 0}
\definecolor{codeidentifier}{RGB}{0, 0, 0}
\definecolor{codelinenumbers}{RGB}{90, 90, 90}

\definecolor{codegreen}{rgb}{0,0.6,0}
\definecolor{codegray}{rgb}{0.5,0.5,0.5}
\definecolor{codepurple}{rgb}{0.58,0,0.82}
\definecolor{backcolour}{rgb}{0.95,0.95,0.92}

\definecolor{aivnGreen}{RGB}{52,79,30}
\definecolor{aivnCodeOutputColor}{RGB}{11,83,148}
\definecolor{captiongray}{gray}{0.15}
\definecolor{captionbg}{RGB}{250,250,250}

\definecolor{foldercolor}{RGB}{0, 102, 204}
\definecolor{filecolor}{RGB}{104, 104, 104}
\definecolor{scriptcolor}{RGB}{180, 50, 50}
\definecolor{pythoncolor}{RGB}{0, 128, 128}

% --- Extra Brand colors for the table ---
\definecolor{brandIndigo}{HTML}{333D91}
\definecolor{rowLavender}{HTML}{EEECFA}
\setlength{\parskip}{0.3em}

% Table
\newtcolorbox{captionbelowboxaivn}{
    colback=aivnGreen,
    colframe=aivnGreen,
    boxrule=0pt,
    arc=4pt,
    left=4pt,
    right=4pt,
    top=2pt,
    bottom=2pt,
    width=\linewidth,
    boxsep=4pt,
    halign=center,
    % fontupper=\bfseries\color{white}
    fontupper=\color{white}
}

% Figure
\newtcolorbox{figurecardcaptionaivn}{
  colback=aivnGreen,        % Màu nền (background color)
  colframe=aivnGreen,       % Màu viền box (border color)
  boxrule=0pt,              % Độ dày viền box (0pt = không viền)
  arc=2pt,                  % Độ bo góc (2pt là bo nhẹ, tăng để bo tròn hơn)
  top=3pt,                  % Padding phía trên nội dung
  bottom=3pt,               % Padding phía dưới nội dung
  left=10pt,                % Padding bên trái
  right=10pt,               % Padding bên phải
  width=\linewidth,         % Chiều rộng bằng toàn bộ dòng
  fontupper=\small\fontsize{11pt}{13pt}\selectfont\color{white}
}

% Tạo counter riêng
\newcounter{hinh}
% Định nghĩa macro caption Hình
\newcommand{\capfig}[1]{%
  \refstepcounter{hinh}%
  % \textbf{Hình~\thehinh:}~#1
    Hình~\thehinh:~#1
}

\newtcolorbox{figurecaptionmodern}{
  colback=captionbg,
  colframe=aivnGreen,
  boxrule=1.5pt,
  bottomrule=2pt,
  toptitle=0pt,
  bottomtitle=0pt,
  left=6pt,
  right=6pt,
  top=4pt,
  bottom=2pt,
  width=\linewidth,
  fontupper=\small\fontsize{11pt}{13pt}\selectfont\color{black},
  halign=center
}

% Code Block -----------------------------------------------------------------
\newtcolorbox{aivncodebox}[1][]{%
  enhanced,
  colback=gray!10,
  colframe=aivnGreen,
  coltitle=white,
  colbacktitle=aivnGreen,
  title=#1,
  fonttitle=\bfseries\sffamily,
  listing only,
  breakable,
  listing engine=listings,
  listing options={
    language=Python,
    basicstyle=\ttfamily\footnotesize,
    commentstyle=\color{ForestGreen},
    keywordstyle=\color{blue},
    stringstyle=\color{red!70!black},
    numbers=left, upquote=true,
    numberstyle=\tiny\color{gray},
    numbersep=5pt,
    breaklines=true,
    showstringspaces=false,
    xleftmargin=1.5em,
    framexleftmargin=1.5em
  }
}

% Code Output Block -----------------------------------------------------------------
\newtcolorbox{aivnCodeOutput}[1][]{%
  enhanced,
  colback=gray!10,
  colframe=aivnCodeOutputColor,
  coltitle=white,
  colbacktitle=aivnCodeOutputColor,
  title=#1,
  fonttitle=\bfseries\sffamily,
  listing only,
  breakable,
  listing engine=listings,
  listing options={
    language=Python,
    basicstyle=\ttfamily\footnotesize,
    commentstyle=\color{ForestGreen},
    keywordstyle=\color{blue},
    stringstyle=\color{red!70!black},
    numbers=left,
    numberstyle=\tiny\color{gray},
    numbersep=5pt,
    breaklines=true,
    showstringspaces=false,
    xleftmargin=1.5em,
    framexleftmargin=1.5em
  }
}

% Last argument allows more tcolorbox options to be added

% --- Guideline/Tip box ---
\newtcolorbox{guidelinebox}[2][]{
    enhanced, colback=yellow!10!white, colframe=orange!90!black, boxrule=1pt, arc=4pt,
    fonttitle=\bfseries, attach boxed title to top left={yshift=-2mm, xshift=3mm},
    colbacktitle=orange!90!black, coltitle=white,
    title=\faLightbulb\hspace{1mm} #2, #1
}

\newtcolorbox{notebox}[1][]{
  enhanced,  colback=yellow!10, colframe=orange!80, boxrule=1pt,
  arc=4pt, fontupper=\small, attach boxed title to top left={yshift=-2mm, xshift=3mm},
  colbacktitle=orange!80,  coltitle=white, fonttitle=\bfseries,
  title=\faInfoCircle~\textbf{Lưu ý},  #1
}

\newtcolorbox{infobox}[1]{%
    enhanced,
    breakable,
    colback=green!3,
    colframe=green!3,
    arc=5pt,
    borderline west={3pt}{0pt}{aivnGreen},
    attach boxed title to top left={
        yshift*=-\tcboxedtitleheight/2,
        xshift=-0.1em
    },
    colbacktitle=aivnGreen, %
    coltitle=white,
    fonttitle=\bfseries,
    boxed title style={
        arc=4pt,
        top=2pt,
        bottom=2pt,
        left=8pt,
        right=8pt,
    },
    title={\faInfoCircle\quad #1}
}
\newtcolorbox{examplebox}[1]{%
    enhanced,
    breakable,
    colback=brandIndigo!5,
    colframe=brandIndigo!5,
    arc=5pt,
    borderline west={3pt}{0pt}{brandIndigo},
    attach boxed title to top left={
        yshift*=-\tcboxedtitleheight/2,
        xshift=-0.1em
    },
    colbacktitle=brandIndigo,
    coltitle=white,
    fonttitle=\bfseries,
    boxed title style={
        arc=4pt,
        top=2pt,
        bottom=2pt,
        left=8pt,
        right=8pt,
    },
    title={\faLightbulb\quad #1}
}
% Prompt Block ---------------------------------------------------------------
\definecolor{promptBlue}{HTML}{2F5D8C}
\definecolor{promptBg}{HTML}{F6F9FC}
\definecolor{promptBorder}{HTML}{B7C9DC}

% Dùng cho prompt đã viết theo cú pháp LaTeX bình thường
\newtcolorbox{promptbox}[1][Prompt]{%
    enhanced,
    breakable,
    colback=promptBg,
    colframe=promptBorder,
    coltitle=white,
    colbacktitle=promptBlue,
    fonttitle=\bfseries\sffamily,
    title={#1},
    boxrule=0.8pt,
    arc=4pt,
    left=8pt,
    right=8pt,
    top=2pt,
    bottom=1pt,
    before skip=8pt,
    after skip=8pt,
    borderline west={3pt}{0pt}{promptBlue},
    fontupper=\small\sffamily,
}

% Dùng cho prompt raw: giữ nguyên xuống dòng, không highlight keyword
\newtcblisting{promptboxraw}[1][Prompt]{%
    enhanced,
    breakable,
    listing only,
    listing engine=listings,
    colback=promptBg,
    colframe=promptBorder,
    coltitle=white,
    colbacktitle=promptBlue,
    fonttitle=\bfseries\sffamily,
    title={#1},
    boxrule=0.8pt,
    arc=4pt,
    left=8pt,
    right=8pt,
    top=-4pt,
    bottom=-4pt,
    before skip=4pt,
    after skip=4pt,
    borderline west={3pt}{0pt}{promptBlue},
    listing options={
        language={},
        basicstyle=\small\ttfamily,
        columns=fullflexible,
        keepspaces=true,
        breaklines=true,
        showstringspaces=false,
        numbers=none,
        frame=none,
        keywordstyle=\color{black},
        commentstyle=\color{black},
        stringstyle=\color{black},
        identifierstyle=\color{black}
    }
}


%--------------------------------------------------------------------
\lstdefinestyle{mystyle}{
    backgroundcolor=\color{codebackground},
    commentstyle=\color{codecomment},
    keywordstyle=\color{codekeyword},
    stringstyle=\color{codestring},
    identifierstyle=\color{codeidentifier},
    basicstyle=\ttfamily\footnotesize,
    breakatwhitespace=false,
    breaklines=true,
    captionpos=b,
    keepspaces=true,
    numbers=left,
    numberstyle=\tiny\color{codelinenumbers},
    numbersep=5pt,
    showspaces=false,
    showstringspaces=false,
    showtabs=false,
    tabsize=2
}

\lstset{style=mystyle}

\hypersetup{
    colorlinks=true,
    linkcolor=black,
    filecolor=magenta,
    citecolor=green,
    urlcolor=blue,
    pdftitle={Nguyên nhân và cơ chế của Vanishing Gradient},
    pdfauthor={Tâm, Trang, Trí, Khanh, Khoa, Long},
    pdflang={vi},
    pdfpagemode=FullScreen,
}

%--------------------------------------------------------------------
\newcounter{choice}
\renewcommand\thechoice{($\alph{choice}$)}
\newcommand\choicelabel{\thechoice}

%--------------------------------------------------------------------
\newenvironment{choices}%
  {\list{\choicelabel}%
     {\usecounter{choice}\def\makelabel##1{\hss\llap{##1}}%
       \settowidth{\leftmargin}{W.\hskip\labelsep\hskip 0 em}%
       \def\choice{%
         \item
       } % choice
       \labelwidth\leftmargin\advance\labelwidth-\labelsep
       \topsep=0pt
       \partopsep=0pt
     }%
  }%
  {\endlist}

%--------------------------------------------------------------------
\def\CChoice{%
      \choice
      \addanswer{\theenumi}{\thechoice}%
    }
    \let\CChoice\CChoice
%    \par % Uncomment this to have choices always start a new line
   % \let\par\@empty
    % If we're continuing the textbf containing the question,
    % then leave a bit of space before the first choice:
    \ifvmode\else\enskip\fi
    \ignorespaces
  %
  {}
\makeatother
\newbox\allanswers
\setbox\allanswers=\hbox{}
\newcommand{\addanswer}[2]{%
  \global\setbox\allanswers=\hbox{\unhbox\allanswers \quad #1.~#2}%
}
\newcommand{\showanswers}{%
  \vfill
  \begin{center}
    Answers
  \end{center}
  \unhbox\allanswers
}

\makeatletter
\newenvironment{multichoices}[1][1]{%
\begin{multicols}{#1}}{%
\end{multicols}}
\makeatother

\pagestyle{fancy}

\DeclareMathOperator*{\argmax}{arg\,max}
\DeclareMathOperator*{\argmin}{arg\,min}
\DeclareMathOperator{\E}{\mathbb{E}}

% Outline --------------------------------------------------------------------
% Section numbering ----------------------------------------------------------
\renewcommand{\thesection}{\Roman{section}.}
\renewcommand{\thesubsection}{\thesection\arabic{subsection}.}
\renewcommand{\thesubsubsection}{\thesubsection\arabic{subsubsection}.}

% Paper-like section title sizes --------------------------------------------
% Title paper: 17pt
% Section should be clearly smaller than title but still prominent.
\titleformat{\section}
  {\normalfont\fontsize{16pt}{20pt}\selectfont\bfseries}
  {\thesection}
  {0.85em}
  {}

\titleformat{\subsection}
  {\normalfont\fontsize{14pt}{16pt}\selectfont\bfseries}
  {\thesubsection}
  {0.75em}
  {}

\titleformat{\subsubsection}
  {\normalfont\fontsize{12pt}{15pt}\selectfont\bfseries}
  {\thesubsubsection}
  {0.65em}
  {}

% Paper-like compact spacing -------------------------------------------------
\titlespacing*{\section}
  {0pt}
  {1.25em plus 0.2em minus 0.1em}
  {0.65em plus 0.1em minus 0.1em}

\titlespacing*{\subsection}
  {0pt}
  {1.0em plus 0.2em minus 0.1em}
  {0.45em plus 0.1em minus 0.1em}

\titlespacing*{\subsubsection}
  {0pt}
  {0.85em plus 0.15em minus 0.1em}
  {0.35em plus 0.1em minus 0.1em}


% Logo + Title ---------------------------------------------------------------

\setlength{\droptitle}{-7em}

% Title horizontal rule
\newcommand{\TitleRule}{%
  \rule{0.90\textwidth}{0.6pt}
}

% Chiều cao vùng title giữa 2 rule
\newlength{\TitleBlockHeight}
\setlength{\TitleBlockHeight}{2.2cm}

\pretitle{%


  % Logo row
  \begin{tabular}{@{}c@{\hspace{1.0em}}c@{\hspace{1.0em}}c@{}}
    \raisebox{-0.35\height}{\includegraphics[height=2.0cm]{assets/aio2026_logo.pdf}}
  \end{tabular}

  %\vspace{0.2em}

\begin{center}
  \TitleRule
  \par\nointerlineskip
}

\title{%
  \parbox[c][\TitleBlockHeight][c]{0.90\textwidth}{%
    \centering
    {\fontsize{20pt}{24pt}\selectfont\bfseries Nguyên nhân và cơ chế của Vanishing Gradient}
  }%
}

\posttitle{%
  \par\nointerlineskip
  \TitleRule
  \end{center}
}

\preauthor{%
  \begin{center}
  \fontsize{11pt}{13pt}\selectfont
}

\author{%
  \textbf{Tâm}\hspace{1.6em}
  \textbf{Trang}\hspace{1.6em}
  \textbf{Trí}\hspace{1.6em}
  \textbf{Khanh}\hspace{1.6em}
  \textbf{Khoa}\hspace{1.6em}
  \textbf{Long}\\
  \vspace{1em}
  \textbf{Proofreading: Trang, Trí}
}

\postauthor{%
  \par
  \end{center}
}

\date{}
% Header --------------------------------------------------------------------
\pagestyle{fancy}
\fancyhf{}
\lhead{\href{https://www.facebook.com/aivietnam.edu.vn}{\textcolor{aivnGreen}{\bfseries AI VIET NAM (AIO2026)}}}
\rhead{\href{https://aivietnam.edu.vn/}{\textcolor{aivnGreen}{\bfseries aivietnam.edu.vn}}}
\fancyfoot[r]{\textcolor{aivnGreen}{\bfseries Trang \thepage}}
\fancyfoot[l]{%
  \href{https://zalo.me/84911118758}{%
    \textcolor{aivnGreen}{\bfseries SĐT/Zalo: 0911118758 (Dr. Hồng Phúc)}%
  }%
}
\renewcommand{\headrulewidth}{1.0pt}
\renewcommand{\headrule}{{\color{aivnGreen}\hrule width\headwidth height\headrulewidth \vskip-\headrulewidth}}

\renewcommand{\footrulewidth}{1.0pt}
\renewcommand{\footrule}{{\color{aivnGreen}\hrule width\headwidth height\headrulewidth \vskip-\headrulewidth}}


\fancypagestyle{firstpage}{
    \fancyhf{}
    \fancyfoot[r]{\textcolor{aivnGreen}{\bfseries Trang \thepage}}
    \fancyfoot[l]{%
      \href{https://zalo.me/84911118758}{%
        \textcolor{aivnGreen}{\bfseries SĐT/Zalo: 0911118758 (Dr. Hồng Phúc)}%
      }%
    }
    \renewcommand{\headrulewidth}{0pt}
    \renewcommand{\footrulewidth}{1.0pt}
    \renewcommand{\footrule}{%
        {\color{aivnGreen}\hrule width\headwidth height\footrulewidth}
    }
}


\begin{document}
\maketitle
\thispagestyle{firstpage}
\vspace{-3em}

\tableofcontents
\clearpage

\section{Giới thiệu về Vanishing Gradient}

Trong deep learning nói riêng và học máy nói chung, sự biến mất đạo hàm
(Vanishing Gradient) là một trong những hiện tượng gây ra nhiều khó khăn trong
quá trình huấn luyện các mạng nơ-ron sâu. Hiện tượng này thường xảy ra khi
gradient phải lan truyền ngược qua quá nhiều lớp, hoặc qua một chuỗi tuần tự
dài trong mạng hồi quy (Recurrent Neural Network). Khi gradient truyền đến các
lớp đầu tiên, giá trị cập nhật tham số có thể trở nên cực kỳ nhỏ, khiến các lớp
này gần như không được học trong suốt quá trình huấn luyện. Vì vậy, mô hình khó
nắm bắt các biểu diễn phức tạp của dữ liệu, từ đó bị giới hạn về độ sâu và hiệu
năng.


\section{Kiến trúc nền tảng của mạng Neural Network}
\textit{Người phụ trách: Trang, Trí}
 
Mạng neural nhân tạo là một hàm tham số hóa $f_{\theta}:\mathbb{R}^{n}\to\mathbb{R}^{K}$ được xây dựng bằng cách xếp chồng nhiều phép biến đổi đơn giản. Mỗi phép biến đổi gồm một bước \emph{tuyến tính} và một bước \emph{phi tuyến}; chính sự kết hợp này cho phép mạng xấp xỉ những quan hệ rất phức tạp giữa đầu vào và đầu ra. Phần này trình bày các khối cơ bản, cách chúng được ghép lại và cách mạng học từ dữ liệu, làm nền tảng để giải thích hiện tượng vanishing gradient ở phần sau.
 
\subsection{Neuron, trọng số và bias}
Neuron là đơn vị tính toán nhỏ nhất của mạng. Với vector đầu vào $\mathbf{x}\in\mathbb{R}^{n}$, một lớp gồm $m$ neuron trước hết thực hiện phép biến đổi tuyến tính
\begin{equation}
    \mathbf{z}=\mathbf{W}\mathbf{x}+\mathbf{b},
    \qquad \mathbf{W}\in\mathbb{R}^{m\times n},\quad \mathbf{b}\in\mathbb{R}^{m},
    \label{eq:linear}
\end{equation}
rồi đi qua hàm kích hoạt phi tuyến $\sigma$ áp dụng theo từng phần tử để cho đầu ra $\mathbf{a}=\sigma(\mathbf{z})$. Với neuron thứ $j$, ta có $z_j=\sum_{i=1}^{n}W_{ji}x_i+b_j$, tức là tổng có trọng số của các đặc trưng đầu vào cộng thêm một hằng số.
 
\begin{itemize}
    \item \textbf{Trọng số $\mathbf{W}$} cho biết mức độ và chiều ảnh hưởng của từng đặc trưng lên neuron: trọng số dương khuếch đại, trọng số âm ức chế, trọng số gần $0$ khiến đặc trưng gần như bị bỏ qua. Đây là nhóm tham số chính được điều chỉnh trong quá trình huấn luyện.
    \item \textbf{Bias $\mathbf{b}$} dịch chuyển ngưỡng kích hoạt của neuron, giúp mô hình biểu diễn được các hàm không đi qua gốc tọa độ. Về mặt hình học, $\mathbf{W}$ quyết định hướng của siêu phẳng phân chia, còn $\mathbf{b}$ quyết định vị trí của nó.
    \item \textbf{Hàm kích hoạt $\sigma$} đưa tính phi tuyến vào mạng. Nếu bỏ $\sigma$, tích của nhiều ma trận vẫn là một ma trận, nên mạng sâu chỉ tương đương một phép biến đổi tuyến tính duy nhất.
\end{itemize}
 
Ba hàm kích hoạt thường gặp, cùng đạo hàm của chúng, được tóm tắt trong Bảng~\ref{tab:activation}. Cột cuối sẽ đặc biệt quan trọng ở Mục~\ref{sec:vanishing}.
 
\begin{table}[htbp]
    \centering
    \caption{Các hàm kích hoạt phổ biến và độ lớn đạo hàm.}
    \label{tab:activation}
    \begin{tabular}{@{}llcc@{}}
        \toprule
        Hàm & Công thức & Miền giá trị & $\max\sigma'(z)$ \\
        \midrule
        Sigmoid & $\dfrac{1}{1+e^{-z}}$ & $(0,1)$ & $0{,}25$ \\[6pt]
        Tanh & $\dfrac{e^{z}-e^{-z}}{e^{z}+e^{-z}}$ & $(-1,1)$ & $1$ \\[6pt]
        ReLU & $\max(0,z)$ & $[0,\infty)$ & $1$ \\
        \bottomrule
    \end{tabular}
\end{table}
 
\subsection{Các lớp của mạng neural}
Mạng truyền thẳng (feedforward) được tổ chức thành ba loại lớp:
\begin{itemize}
    \item \textbf{Lớp đầu vào (input layer)} nhận dữ liệu thô $\mathbf{x}$ như điểm ảnh, đặc trưng dạng bảng hay embedding của từ. Lớp này không tính toán mà chỉ đưa dữ liệu vào mạng.
    \item \textbf{Các lớp ẩn (hidden layers)} biến đổi dữ liệu qua nhiều bước liên tiếp. Thông thường, các lớp đầu học đặc trưng đơn giản (cạnh, tần suất từ), còn các lớp sau kết hợp chúng thành biểu diễn trừu tượng hơn (hình dạng, ngữ nghĩa).
    \item \textbf{Lớp đầu ra (output layer)} tạo dự đoán $\hat{\mathbf{y}}$. Hàm kích hoạt của lớp này phụ thuộc bài toán: tuyến tính cho hồi quy, sigmoid cho phân loại nhị phân, softmax cho phân loại nhiều lớp.
\end{itemize}
 
Thông tin đi từ lớp này sang lớp kế tiếp. Với mạng gồm $L$ lớp (không tính lớp đầu vào) và quy ước $\mathbf{a}^{(0)}=\mathbf{x}$, ta có
\begin{equation}
    \mathbf{z}^{(l)}=\mathbf{W}^{(l)}\mathbf{a}^{(l-1)}+\mathbf{b}^{(l)},
    \qquad
    \mathbf{a}^{(l)}=\sigma^{(l)}\!\bigl(\mathbf{z}^{(l)}\bigr),
    \qquad l=1,\dots,L.
    \label{eq:layer}
\end{equation}
Tập tham số của toàn mạng là $\theta=\{\mathbf{W}^{(l)},\mathbf{b}^{(l)}\}_{l=1}^{L}$.
 
\subsection{Forward pass và hàm mất mát}
\textbf{Forward pass} là việc áp dụng công thức~\eqref{eq:layer} lần lượt từ lớp $1$ đến lớp $L$ để thu được dự đoán $\hat{\mathbf{y}}=\mathbf{a}^{(L)}$. Trong quá trình này, các giá trị trung gian $\mathbf{z}^{(l)}$ và $\mathbf{a}^{(l)}$ được lưu lại vì bước lan truyền ngược sẽ cần đến chúng.
 
Sau đó, \textbf{hàm mất mát} (loss) đo độ lệch giữa dự đoán $\hat{\mathbf{y}}$ và nhãn thật $\mathbf{y}$. Hai lựa chọn phổ biến là
\begin{align}
    \mathcal{L}_{\mathrm{MSE}} &= \frac{1}{N}\sum_{i=1}^{N}\bigl\lVert \hat{\mathbf{y}}_i-\mathbf{y}_i\bigr\rVert_2^{2}
    && \text{(hồi quy)}, \\
    \mathcal{L}_{\mathrm{CE}} &= -\frac{1}{N}\sum_{i=1}^{N}\sum_{k=1}^{K} y_{ik}\log\hat{y}_{ik}
    && \text{(phân loại)}.
\end{align}
Loss là một đại lượng vô hướng: nó cho biết mô hình đang sai bao nhiêu và là mục tiêu cần cực tiểu hóa. Đồng thời, nó là điểm xuất phát của lan truyền ngược.
 
\subsection{Backpropagation và cập nhật tham số}
Để giảm loss, ta cần biết mỗi tham số ảnh hưởng đến loss như thế nào, tức là \textbf{gradient} $\partial\mathcal{L}/\partial\mathbf{W}^{(l)}$ và $\partial\mathcal{L}/\partial\mathbf{b}^{(l)}$. \textbf{Backpropagation} tính đồng thời gradient của mọi lớp một cách hiệu quả nhờ \textbf{quy tắc chuỗi} (chain rule), đi ngược từ lớp đầu ra về lớp đầu vào và tái sử dụng kết quả trung gian.
 
Đặt sai số cục bộ $\boldsymbol{\delta}^{(l)}=\partial\mathcal{L}/\partial\mathbf{z}^{(l)}$. Khi đó
\begin{align}
    \boldsymbol{\delta}^{(L)} &= \frac{\partial\mathcal{L}}{\partial\mathbf{a}^{(L)}}\odot\sigma'\!\bigl(\mathbf{z}^{(L)}\bigr), \label{eq:delta-out}\\
    \boldsymbol{\delta}^{(l)} &= \Bigl(\mathbf{W}^{(l+1)}\Bigr)^{\!\top}\boldsymbol{\delta}^{(l+1)}\odot\sigma'\!\bigl(\mathbf{z}^{(l)}\bigr), \label{eq:delta-rec}\\
    \frac{\partial\mathcal{L}}{\partial\mathbf{W}^{(l)}} &= \boldsymbol{\delta}^{(l)}\Bigl(\mathbf{a}^{(l-1)}\Bigr)^{\!\top},
    \qquad
    \frac{\partial\mathcal{L}}{\partial\mathbf{b}^{(l)}} = \boldsymbol{\delta}^{(l)}, \label{eq:grad}
\end{align}
với $\odot$ là tích Hadamard (nhân từng phần tử). Có gradient rồi, tham số được cập nhật theo \textbf{gradient descent} với tốc độ học $\eta>0$:
\begin{equation}
    \mathbf{W}^{(l)}\leftarrow\mathbf{W}^{(l)}-\eta\,\frac{\partial\mathcal{L}}{\partial\mathbf{W}^{(l)}},
    \qquad
    \mathbf{b}^{(l)}\leftarrow\mathbf{b}^{(l)}-\eta\,\frac{\partial\mathcal{L}}{\partial\mathbf{b}^{(l)}}.
    \label{eq:gd}
\end{equation}
Toàn bộ một vòng huấn luyện gồm bốn bước lặp lại: \emph{forward} $\rightarrow$ \emph{tính loss} $\rightarrow$ \emph{backward} $\rightarrow$ \emph{cập nhật tham số}.
 
\begin{notebox}
Công thức đệ quy~\eqref{eq:delta-rec} cho thấy $\boldsymbol{\delta}^{(l)}$ được tạo ra bằng cách nhân liên tiếp các ma trận $(\mathbf{W}^{(k)})^{\top}$ và các đạo hàm $\sigma'$ qua từng lớp. Một chuỗi tích dài như vậy có thể co lại hoặc phình ra rất nhanh. Đây chính là gốc rễ của hiện tượng vanishing gradient được phân tích ở phần tiếp theo.
\end{notebox}
 
 
\section{Dấu hiệu nhận biết Vanishing Gradient}
\label{sec:vanishing}
\textit{Người phụ trách: Trang, Trí}
 
\subsection{Cơ chế hình thành}
Khai triển đệ quy~\eqref{eq:delta-rec} cho thấy sai số tại lớp $l$ là tích của nhiều thừa số:
\begin{equation}
    \boldsymbol{\delta}^{(l)}
    =\Biggl[\prod_{k=l+1}^{L}\operatorname{diag}\!\bigl(\sigma'(\mathbf{z}^{(k-1)})\bigr)\bigl(\mathbf{W}^{(k)}\bigr)^{\!\top}\Biggr]\boldsymbol{\delta}^{(L)}.
    \label{eq:product}
\end{equation}
Nếu chuẩn của mỗi thừa số nhỏ hơn $1$, gradient co lại gần như theo hàm mũ khi đi ngược về các lớp đầu. Ví dụ, với sigmoid ta có $\sigma'(z)\le 0{,}25$; chỉ riêng thành phần này, sau $10$ lớp gradient đã bị nhân với không quá $0{,}25^{10}\approx 10^{-6}$. Khi đó các lớp gần đầu vào nhận được tín hiệu học gần như bằng không, dù lớp cuối vẫn học bình thường. Hiện tượng này được gọi là \textbf{vanishing gradient}.
 
\subsection{Các dấu hiệu quan sát được}
Các dấu hiệu dưới đây có thể được kiểm tra trực tiếp trong quá trình huấn luyện; Bảng~\ref{tab:signs} tóm tắt cách phát hiện từng dấu hiệu.
 
\begin{enumerate}
    \item \textbf{Gradient ở các lớp gần đầu vào rất nhỏ hoặc gần bằng $0$.} Chuẩn gradient $\lVert\nabla_{\mathbf{W}^{(l)}}\mathcal{L}\rVert$ của lớp đầu nhỏ hơn lớp cuối nhiều bậc độ lớn (chẳng hạn $10^{-8}$ so với $10^{-2}$), và giảm đều theo chiều từ lớp cuối về lớp đầu.
    \item \textbf{Trọng số của các lớp đầu thay đổi rất ít qua nhiều epoch.} Bước cập nhật $\eta\,\nabla\mathcal{L}$ gần như bằng $0$ nên các trọng số này gần như đứng yên so với giá trị khởi tạo.
    \item \textbf{Loss giảm rất chậm, sớm tạo plateau hoặc mô hình không hội tụ như kỳ vọng.} Learning curve gần như phẳng dù vẫn còn nhiều epoch và không có dấu hiệu overfitting (loss huấn luyện cũng không giảm).
    \item \textbf{Các activation rơi vào vùng bão hòa của sigmoid/tanh.} Khi $|z|$ lớn, đầu ra tiến sát $0$ hoặc $1$ (sigmoid), $\pm1$ (tanh), khiến $\sigma'(z)\approx 0$ và chặn dòng gradient đi qua. Histogram activation khi đó dồn về hai đầu mút.
    \item \textbf{Lớp cuối học được trong khi các lớp đầu gần như không học.} Mô hình chỉ tinh chỉnh phần đầu ra trên những đặc trưng gần như ngẫu nhiên, nên độ chính xác bị chặn ở mức thấp.
    \item \textbf{Với RNN, mô hình ghi nhớ kém các phụ thuộc dài hạn.} Gradient lan truyền qua nhiều bước thời gian bị co lại, nên mô hình chỉ nắm được ngữ cảnh gần và bỏ sót thông tin ở xa.
\end{enumerate}
 
\begin{table}[htbp]
    \centering
    \caption{Tóm tắt dấu hiệu vanishing gradient và cách kiểm tra.}
    \label{tab:signs}
    \renewcommand{\arraystretch}{1.25}
    \begin{tabularx}{\textwidth}{@{}>{\raggedright\arraybackslash}p{0.30\textwidth} >{\raggedright\arraybackslash}X >{\raggedright\arraybackslash}X@{}}
        \toprule
        Dấu hiệu & Đại lượng cần theo dõi & Biểu hiện điển hình \\
        \midrule
        Gradient lớp đầu $\approx 0$ & Gradient norm theo từng lớp & Giảm theo cấp số nhân về phía lớp đầu \\
        Trọng số lớp đầu ít đổi & $\lVert\mathbf{W}^{(l)}_t-\mathbf{W}^{(l)}_0\rVert$ theo epoch & Gần như phẳng \\
        Loss giảm chậm, plateau & Learning curve (train loss) & Đường cong phẳng sớm \\
        Activation bão hòa & Histogram của $\mathbf{a}^{(l)}$ & Dồn về $0/1$ hoặc $\pm1$ \\
        Chỉ lớp cuối học & So sánh tốc độ thay đổi giữa các lớp & Lớp cuối thay đổi, lớp đầu đứng yên \\
        RNN quên ngữ cảnh xa & Gradient theo bước thời gian & Suy giảm nhanh khi lùi xa \\
        \bottomrule
    \end{tabularx}
\end{table}
 
\subsection{Phân biệt với các nguyên nhân khác}
Loss giảm chậm chưa chắc do vanishing gradient. Cần loại trừ các nguyên nhân thường gặp khác như tốc độ học $\eta$ quá nhỏ, lỗi cài đặt (quên gọi \texttt{zero\_grad}, sai nhãn), hay ``dead ReLU'' (neuron luôn xuất ra $0$). Kết luận chỉ đáng tin khi \emph{đồng thời} thấy gradient các lớp đầu nhỏ hơn hẳn các lớp cuối \emph{và} activation bão hòa (hoặc đạo hàm nhỏ) ở những lớp đó.

\section{Cơ chế Vanishing Gradient trong Backpropagation}
\textit{Người phụ trách: Khanh}

\subsection{Chain rule trong mạng nhiều lớp}
Với mạng sâu, gradient tại một lớp sớm là tích của nhiều đạo hàm cục bộ:
\[
\frac{\partial \mathcal{L}}{\partial \mathbf{h}^{(l)}}
=
\left(\prod_{k=l+1}^{L}
\frac{\partial \mathbf{h}^{(k)}}{\partial \mathbf{h}^{(k-1)}}\right)
\frac{\partial \mathcal{L}}{\partial \mathbf{h}^{(L)}}.
\]

\subsection{Sự co nhỏ qua phép nhân lặp}
Nếu chuẩn của các Jacobian thường nhỏ hơn $1$, tích lặp qua nhiều lớp sẽ giảm theo cấp số nhân, khiến gradient truyền về các lớp đầu trở nên rất nhỏ.

\subsection{Liên hệ giữa đạo hàm activation và trọng số}
Phân tích cách $f'(z)$ và phổ của ma trận trọng số cùng quyết định độ lớn gradient.

\subsection{Ví dụ số tối giản}
Đề xuất tính tay một chuỗi gồm nhiều đạo hàm bằng $0.25$ để cho thấy $0.25^L$ nhanh chóng tiến về $0$ khi $L$ tăng.

\section{Các yếu tố gây Vanishing Gradient}

\subsection{Hàm activation}
\textit{Người phụ trách: Khanh}

\begin{itemize}
    \item Sigmoid và tanh có vùng bão hòa với đạo hàm rất nhỏ.
    \item Sigmoid có đạo hàm cực đại chỉ bằng $0.25$.
    \item So sánh với ReLU và các biến thể để làm rõ ảnh hưởng của activation.
\end{itemize}

\subsection{Khởi tạo và độ lớn của weight}
\textit{Người phụ trách: Khoa}

\begin{itemize}
    \item Trọng số quá nhỏ làm cả tín hiệu và gradient co lại qua từng lớp.
    \item Trọng số quá lớn có thể đẩy activation vào vùng bão hòa.
    \item Liên hệ với Xavier/Glorot initialization và He initialization.
\end{itemize}

\subsection{Độ sâu của mạng}
\textit{Người phụ trách: Khoa}

\begin{itemize}
    \item Số lượng thừa số trong chain rule tăng theo độ sâu.
    \item Gradient suy giảm theo cấp số nhân nếu các thừa số có chuẩn nhỏ hơn $1$.
    \item Các lớp gần đầu vào chịu ảnh hưởng nghiêm trọng nhất.
\end{itemize}

\subsection{Sự suy giảm tín hiệu trong Forward Pass}
\textit{Người phụ trách: Long}

\begin{itemize}
    \item Activation co nhỏ hoặc mất phương sai khi đi qua nhiều lớp.
    \item Tín hiệu đầu vào yếu dần làm biểu diễn ở lớp sâu kém phân biệt.
    \item Liên hệ giữa ổn định forward signal và ổn định backward gradient.
\end{itemize}

\subsection{Vanishing Gradient trong RNN}
\textit{Người phụ trách: Long}

\begin{itemize}
    \item Backpropagation Through Time tạo tích lặp của recurrent weight và đạo hàm activation.
    \item Gradient từ thời điểm xa suy giảm trước khi quay về các bước đầu chuỗi.
    \item Hệ quả là khó học phụ thuộc dài hạn và dễ quên thông tin cũ.
\end{itemize}

\subsection{Ảnh hưởng của kiến trúc mạng}
\textit{Người phụ trách: Tâm}

\begin{itemize}
    \item Mạng tuần tự thuần túy buộc gradient đi qua toàn bộ chuỗi phép biến đổi.
    \item Thiếu skip/residual connection làm đường truyền gradient dài hơn.
    \item Bottleneck quá hẹp hoặc cấu trúc không chuẩn hóa có thể làm tín hiệu suy yếu.
\end{itemize}

\section{Ví dụ minh họa cho các trường hợp Vanishing Gradient}
\textit{Người phụ trách chung: Khanh, Tâm, Long, Khoa}

\subsection{MLP sâu dùng sigmoid}
\textit{Phụ trách: Khanh}

So sánh gradient norm giữa lớp đầu và lớp cuối khi tăng số lớp của một MLP dùng sigmoid.

\subsection{Ảnh hưởng của khởi tạo weight và độ sâu}
\textit{Phụ trách: Khoa}

Giữ nguyên dữ liệu và kiến trúc cơ bản, thay đổi độ sâu cùng phương pháp khởi tạo để quan sát loss và gradient.

\subsection{RNN học phụ thuộc dài hạn}
\textit{Phụ trách: Long}

Dùng một bài toán chuỗi đơn giản, thay đổi sequence length và đo gradient theo khoảng cách thời gian.

\subsection{Mạng có và không có residual connection}
\textit{Phụ trách: Tâm}

So sánh đường truyền gradient của plain network với residual network có cùng độ sâu.

\begin{examplebox}{Đầu ra nên có cho mỗi ví dụ}
Mỗi ví dụ nên gồm: mục tiêu, kiến trúc, cấu hình huấn luyện, đoạn code tối giản, biểu đồ gradient/loss và phần nhận xét ngắn gọn.
\end{examplebox}

\section{Các giải pháp phổ biến}
\textit{Người phụ trách chung: Khanh, Tâm, Long, Khoa}

\subsection{Lựa chọn hàm activation phù hợp}
\textit{Phụ trách: Khanh}

ReLU, Leaky ReLU, ELU/GELU và lưu ý về \textit{dying ReLU}.

\subsection{Khởi tạo trọng số và chuẩn hóa}
\textit{Phụ trách: Khoa}

Xavier/Glorot, He initialization, Batch Normalization và Layer Normalization.

\subsection{Residual/skip connection và thiết kế kiến trúc}
\textit{Phụ trách: Tâm}

Giải thích đường tắt cho gradient, residual block và nguyên tắc thiết kế mạng sâu dễ tối ưu.

\subsection{Giải pháp dành cho RNN}
\textit{Phụ trách: Long}

LSTM, GRU, gradient clipping, truncated BPTT và lựa chọn kiến trúc thay thế khi phù hợp.

\subsection{Theo dõi gradient trong quá trình huấn luyện}
Đề xuất ghi log gradient norm, histogram trọng số/gradient và cảnh báo các lớp có gradient gần bằng $0$.

\begin{guidelinebox}{Thông điệp thực hành}
Không có một giải pháp duy nhất cho mọi mô hình. Nên kết hợp activation phù hợp, initialization đúng, normalization, residual connection và theo dõi gradient để chẩn đoán sớm.
\end{guidelinebox}

\section{Kết luận}
\textit{Người phụ trách: Trang, Trí}

\begin{itemize}
    \item Tóm tắt nguyên nhân cốt lõi: gradient là tích của nhiều đạo hàm/Jacobian nhỏ.
    \item Nhấn mạnh hậu quả: các lớp đầu hoặc các thời điểm xa không nhận đủ tín hiệu học.
    \item Tổng hợp dấu hiệu chẩn đoán và nhóm giải pháp phù hợp cho MLP/CNN/RNN.
    \item Nêu thông điệp cuối: thiết kế kiến trúc, khởi tạo và giám sát quá trình huấn luyện phải được xem xét đồng thời.
\end{itemize}

\end{document}
